\chapter{Синхронные схемы, автоматы и память}
\label{ch:sync}

В предыдущих главах мы научились вычислять (сумматоры, компараторы, мультиплексоры)
и запоминать (триггеры, регистры, счётчики). Остался вопрос, ради которого
всё это затевалось: как соединить вычисление и память в одну схему, которой
можно доверять на высокой частоте и которая ведёт себя по сценарию~--- ждёт
события, реагирует, возвращается в исходное положение? Ответ даёт синхронный
принцип и конечные автоматы. А в конце главы мы заглянем в память: она ---
та же комбинационная логика и те же регистры, только в большом количестве.

\begin{buildbox}
Откройте вкладку \emph{Обучение}, книгу «из блоков», главу «Автоматы», урок
«Детектор 101». Соберите автомат, который выдаёт 1, когда на вход пришли
подряд биты 1, 0, 1. Во вкладке \emph{Диаграммы} подайте на вход $x$
последовательность \code{1010110010} и пронаблюдайте регистр состояния и
выход~$y$. Вопрос на время чтения: через сколько тактов после последней
единицы тройки выход поднимается в 1 и почему?
\end{buildbox}

\section{Синхронное проектирование}

Мы уже знаем два «материала»: комбинационную логику (она вычисляет) и
триггеры (они хранят). Как собрать из них большое надёжное устройство?
Почти все современные цифровые схемы строятся по одному принципу ---
\term{синхронному}:

\begin{itemize}
  \item все триггеры тактируются \textbf{одним} тактовым сигналом;
  \item между регистрами находится только комбинационная логика;
  \item результат вычислений записывается в регистры по фронту, когда все
        переходные процессы и «иголки» уже закончились.
\end{itemize}

\begin{figure}[H]
\centering
\begin{tikzpicture}[rg/.style={draw=FpgaGreen, fill=green!8, thick, minimum width=10mm, minimum height=18mm, font=\sffamily\small},
  cl/.style={draw=FpgaBlue, fill=FpgaLight, thick, cloud, cloud puffs=11, aspect=1.8, minimum width=26mm, minimum height=16mm, font=\sffamily\small, align=center}]
  \node[rg] (r1) at (0,0) {RG};
  \node[cl] (c1) at (2.6,0) {логика};
  \node[rg] (r2) at (5.2,0) {RG};
  \node[cl] (c2) at (7.8,0) {логика};
  \node[rg] (r3) at (10.4,0) {RG};
  \draw[bus, -{Stealth}] (-1.2,0) node[left, text=FpgaDark] {входы} -- (r1);
  \draw[bus, -{Stealth}] (r1) -- (c1); \draw[bus, -{Stealth}] (c1) -- (r2);
  \draw[bus, -{Stealth}] (r2) -- (c2); \draw[bus, -{Stealth}] (c2) -- (r3);
  \draw[bus, -{Stealth}] (r3) -- ++(1.2,0) node[right, text=FpgaDark] {выходы};
  \draw[wire, FpgaBlue] (-0.8,-1.6) node[left, text=FpgaDark] {CLK} -- (10.4,-1.6);
  \foreach \r in {r1,r2,r3} { \draw[wire, FpgaBlue] (\r.south) -- (\r.south |- 0,-1.6); \fill[FpgaBlue] (\r.south |- 0,-1.6) circle (1.4pt); }
\end{tikzpicture}
\caption{Синхронная схема: регистры, разделённые комбинационной логикой, и
один общий такт}
\label{fig:sync}
\end{figure}

\subsection{Максимальная тактовая частота}

За один такт сигнал должен пройти от выхода одного регистра через логику
до входа следующего и успеть «установиться» до фронта. Отсюда главное
уравнение синхронной схемы:
\[
  T_\text{CLK} \ \ge\ t_\text{co} + t_\text{логики} + t_\text{su},
  \qquad
  f_\text{max} = \frac{1}{t_\text{co} + t_\text{логики} + t_\text{su}}.
\]

\begin{figure}[H]
\centering
\begin{tikzpicture}[x=1cm, y=1cm]
  \draw[very thick, FpgaBlue] (0,2.2) -- (1,2.2) -- (1.05,2.8) -- (5,2.8) -- (5.05,2.2) -- (9,2.2) -- (9.05,2.8) -- (11,2.8);
  \node[font=\sffamily\small, anchor=east] at (-0.1,2.5) {CLK};
  \fill[FpgaOrange!40] (1.05,1.2) rectangle (2.0,1.8); \node[font=\sffamily\scriptsize] at (1.52,1.5) {$t_\text{co}$};
  \fill[FpgaBlue!25] (2.0,1.2) rectangle (7.6,1.8); \node[font=\sffamily\scriptsize] at (4.8,1.5) {$t_\text{логики}$ (критический путь)};
  \fill[FpgaRed!30] (7.6,1.2) rectangle (8.6,1.8); \node[font=\sffamily\scriptsize] at (8.1,1.5) {$t_\text{su}$};
  \fill[green!25] (8.6,1.2) rectangle (9.05,1.8);
  \node[font=\sffamily\scriptsize, anchor=west] at (9.2,1.5) {запас};
  \draw[dashed, FpgaGray] (1.05,3.0) -- (1.05,0.6); \draw[dashed, FpgaGray] (9.05,3.0) -- (9.05,0.6);
  \draw[<->, thick] (1.05,0.8) -- (9.05,0.8) node[midway, below, font=\sffamily\small] {период $T_\text{CLK}$};
\end{tikzpicture}
\caption{Бюджет времени одного такта. Если сумма задержек больше периода,
в регистр запишется неверное значение}
\label{fig:budget}
\end{figure}

Пример: $t_\text{co} = 0{,}5$~нс, $t_\text{su} = 0{,}5$~нс, а сигнал проходит
через 10 элементов по 0,5~нс. Тогда
$T \ge 0{,}5 + 5 + 0{,}5 = 6$~нс и $f_\text{max} \approx 166$~МГц. Чем
длиннее цепочка логики между регистрами, тем ниже допустимая частота.

\begin{figure}[H]
\centering
\begin{tikzpicture}
\begin{axis}[
  width=0.78\textwidth, height=5.2cm, xlabel={Число уровней логики на критическом пути}, ylabel={$f_\text{max}$, МГц},
  xmin=1, xmax=20, ymin=0, ymax=520,
  label style={font=\sffamily\small}, tick label style={font=\sffamily\small},
  grid=major, grid style={FpgaGray!20}, axis lines*=left,
]
  \addplot[very thick, FpgaBlue, domain=1:20, samples=40] {1000/(1+0.5*x)};
  \addplot[only marks, mark=*, FpgaRed] coordinates {(10,166.7)};
  \node[font=\sffamily\small, text=FpgaRed, anchor=south west] at (axis cs:10,170) {пример: 166 МГц};
\end{axis}
\end{tikzpicture}
\caption{Максимальная частота падает с ростом числа элементов между регистрами
($t_\text{co}+t_\text{su} = 1$~нс, задержка элемента 0,5~нс)}
\end{figure}

\subsection{Конвейер}

Что делать, если логика слишком длинная? Её можно разрезать на части и
поставить между ними дополнительные регистры. Это \term{конвейер}
(pipeline): каждая «станция» делает свою часть работы за один короткий такт,
и хотя результат для одних данных появляется через несколько тактов, новые
данные можно подавать \emph{каждый} такт. Так работает сборочный конвейер на
заводе --- и все современные процессоры.

\begin{figure}[H]
\centering
\begin{tikzpicture}[c/.style={draw, minimum width=11mm, minimum height=6mm, font=\sffamily\scriptsize}]
  \foreach \t in {1,...,6} \node[font=\sffamily\scriptsize] at (\t*1.15+0.9,0.55) {такт \t};
  \foreach \s/\n [count=\r from 0] in {этап 1/A, этап 2/B, этап 3/C} {
    \node[font=\sffamily\small, anchor=east] at (1.3,-\r*0.75) {\s};
  }
  \foreach \d/\col [count=\k from 0] in {1/FpgaBlue!30, 2/FpgaOrange!40, 3/FpgaGreen!35, 4/FpgaPurple!30} {
    \foreach \r in {0,1,2} {
      \pgfmathsetmacro{\x}{(\k+\r+1)*1.15+0.9}
      \node[c, fill=\col] at (\x,-\r*0.75) {$D_{\d}$};
    }
  }
\end{tikzpicture}
\caption{Трёхэтапный конвейер ($D_1,\dots,D_4$ --- порции данных): на такте 3 одновременно обрабатываются три
порции данных, каждая на своём этапе}
\end{figure}

\subsection{Конвейер на Verilog}
\label{sec:sync-pipe-v}

Возьмём вычисление $y=a\cdot b+c$. В первом варианте умножитель и сумматор
стоят в одном такте, во втором между ними поставлен регистр:

\vfile{\codepath/b9_sync/nopipe.v}{code/b9\_sync/nopipe.v}

\vfile{\codepath/b9_sync/pipe2.v}{code/b9\_sync/pipe2.v}

\textbf{Как читается.} В \code{nopipe} одна строка: результат умножения и
сложения записывается в \code{y} за один такт; путь через логику --- самый
длинный, поэтому предельная частота ниже. В \code{pipe2} два этапа: на первом
произведение попадает в регистр \code{prod}, на втором \code{prod + c1}
записывается в \code{y}. Третий регистр \code{c1} нужен, чтобы слагаемое $c$
«ехало» вместе со своей парой $(a,b)$: без него сложение взяло бы $c$ из
\emph{следующей} порции данных. Задержка результата --- 2 такта, но новая
тройка $(a,b,c)$ принимается \emph{каждый} такт.

Тестбенч \code{tb\_pipe.v} подаёт новую тройку на каждом такте
(\code{a=i+1}, \code{b=i+2}, \code{c=i+3}) и проверяет: \code{nopipe}
отвечает на предыдущем такте, \code{pipe2} --- на два такта позже входа. Он
подтверждает главное свойство конвейера: пропускная способность --- один
результат за такт, хотя каждый результат опаздывает на два.

\textbf{Как запустить в Logra.} Вкладка \term{Код}, вставьте \code{pipe2},
Ctrl+Enter, \term{RTL}: два регистра между умножением и сложением. На
\term{Диаграммах} подайте \code{a=3, b=4, c=5} на один такт и найдите,
через сколько тактов \code{y} станет 17. Затем сравните оценки частоты двух
вариантов на вкладке \term{Gowin} (она считает предельную частоту после
размещения и трассировки): у \code{pipe2} путь через логику короче, и оценка
обычно выше.

\begin{diybox}[Конвейер на Verilog: допиши модуль]
\textbf{Задание.} \stars{1} Задержите сигнал на два такта: выход
повторяет вход с задержкой два такта. \stars{2} Добавьте к \code{pipe2}
выход \code{valid}, равный 1, когда на \code{y} уже настоящий результат (в
первые такты после сброса там мусор). \stars{3} Разрежьте вычисление
$y=a\cdot b+c\cdot d$ на три этапа: два умножения, сложение, регистр
выхода.\\
\textbf{Ожидаемый результат:} \code{valid} поднимается через два такта после
первого входа.
\end{diybox}

\textbf{Образец ответа.}
\begin{vcode}
module delay2 (input clk, d, output reg q);
    reg m = 1'b0;
    always @(posedge clk) begin
        m <= d;      // первый такт задержки
        q <= m;      // второй
    end
endmodule

module pipe2_valid (
    input            clk, in_valid,
    input      [3:0] a, b, c,
    output reg [8:0] y,
    output reg       valid
);
    reg [7:0] prod;
    reg [3:0] c1;
    reg       v1;
    always @(posedge clk) begin
        prod <= a * b;  c1 <= c;  v1 <= in_valid;
        y    <= prod + c1;        valid <= v1;   // флаг идёт по конвейеру вместе с данными
    end
endmodule
\end{vcode}

\begin{diybox}[Такт, критический путь и конвейер]
\textbf{Задание.} Увидеть, что цепочка регистров задерживает результат на
целое число тактов, но принимает новые данные каждый такт, и связать это с
формулой $f_\text{max}$.\\
\textbf{В Logra.} Вкладка \emph{Примеры}, проект \code{moving-average}
(среднее четырёх отсчётов: регистры \code{delay1}\ldots\code{delay3} образуют
цепочку).\\
\stars{1} \emph{Повтори.} Задайте \code{SAMPLE\_TICKS=1}, подайте входной
отсчёт 32 и в \emph{Диаграммах} проследите 20 тактов: на сколько тактов
выход отстаёт от входа?\\
\stars{2} \emph{Измени.} Откройте проект \code{fir-filter}, задайте
коэффициенты $A=1$, $B=0$, $C=-1$ и \code{SCALE\_SHIFT=0}. Подайте отсчёт 32:
ожидаемый отклик $32,\,0,\,-32$ в трёх соседних тактах.\\
\stars{3} \emph{Проект.} На бумаге разрежьте логику с $t_\text{логики}=8$~нс
(задача~1) на два этапа и посчитайте новые $T$ и $f_\text{max}$; затем
убедитесь на диаграммах, что новый отсчёт принимается каждый такт.\\
\textbf{Ожидаемый результат.} Задержка результата в тактах равна числу
регистров на пути, а пропускная способность~--- один отсчёт за такт.\\
\textbf{Как проверить.} Вкладка \emph{Диаграммы}: сравните номер такта
входного импульса и такта отклика; во вкладке \emph{Gowin} сравните оценку
предельной частоты проекта с вашим расчётом по формуле главы.
\end{diybox}

\section{Конечные автоматы}

Многие устройства работают по сценарию: «ждать нажатия, потом сделать то, потом
это, при ошибке вернуться в начало». Такое поведение описывают
\term{конечным автоматом} (FSM, finite state machine). Автомат находится в
одном из конечного числа \term{состояний}; по каждому такту, в зависимости от
текущего состояния и входов, он \term{переходит} в следующее состояние.

Схема любого автомата одинакова: регистр хранит номер текущего состояния,
комбинационная логика вычисляет следующее состояние и выходы.

\begin{figure}[H]
\centering
\begin{tikzpicture}[node distance=12mm]
  \node[block, minimum width=28mm, minimum height=14mm] (next) {Логика\\переходов};
  \node[gblock, right=of next, minimum height=14mm] (st) {Регистр\\состояния};
  \node[block, right=of st, minimum width=28mm, minimum height=14mm] (out) {Логика\\выходов};
  \draw[bus, -{Stealth}] (next) -- node[above, lbl] {след.} (st);
  \draw[bus, -{Stealth}] (st) -- node[above, lbl] {текущее} (out);
  \draw[arr] (out.east) -- ++(1,0) node[right, font=\sffamily\small] {выходы};
  \draw[bus, -{Stealth}] ($(st.east)+(0.6,0)$) -- ++(0,-1.3) -| (next.south);
  \draw[arr] ($(next.west)+(-1.2,0)$) node[left, font=\sffamily\small] {входы} -- (next.west);
  \draw[arr, dashed, FpgaPurple] ($(next.west)+(-0.6,0)$) -- ++(0,1.3) -| node[pos=0.25, above, lbl, text=FpgaPurple] {только у автомата Мили} (out.north);
  \draw[wire, FpgaBlue] (st.south) ++(0,0) -- ++(0,-0.5) node[below, font=\sffamily\scriptsize, text=FpgaDark] {CLK};
\end{tikzpicture}
\caption{Структура конечного автомата. У автомата \textbf{Мура} выходы зависят
только от состояния, у автомата \textbf{Мили} --- ещё и от входов (пунктир)}
\label{fig:fsm_struct}
\end{figure}

\subsection{Пример: детектор последовательности}

Построим автомат, который следит за входом $x$ (по одному биту за такт) и
выдаёт $y=1$, когда последние три принятых бита равны \code{101}. Такие
детекторы ищут стартовые метки в потоке данных.

\begin{figure}[H]
\centering
\begin{tikzpicture}[
  st/.style={circle, draw=FpgaDark, very thick, minimum size=17mm, font=\sffamily\small, align=center, fill=FpgaLight},
  tr/.style={-{Stealth[length=3mm]}, thick, FpgaDark},
  lb/.style={font=\sffamily\small, fill=white, inner sep=1pt}]
  \node[st] (s0) at (0,0) {S0\\\scriptsize $y=0$};
  \node[st] (s1) at (3.6,0) {S1\\\scriptsize «1»\\\scriptsize $y=0$};
  \node[st] (s2) at (7.2,0) {S2\\\scriptsize «10»\\\scriptsize $y=0$};
  \node[st, fill=green!20] (s3) at (10.8,0) {S3\\\scriptsize «101»\\\scriptsize $y=1$};
  \draw[tr] (s0) -- node[lb, above] {1} (s1);
  \draw[tr] (s1) -- node[lb, above] {0} (s2);
  \draw[tr] (s2) -- node[lb, above] {1} (s3);
  \draw[tr] (s0) edge[loop above, looseness=6] node[lb, above] {0} (s0);
  \draw[tr] (s1) edge[loop above, looseness=6] node[lb, above] {1} (s1);
  \draw[tr] (s2) to[bend left=35] node[lb, below] {0} (s0);
  \draw[tr] (s3) to[bend left=32] node[lb, below] {0} (s2);
  \draw[tr] (s3) to[bend right=40] node[lb, above] {1} (s1);
  \draw[tr, FpgaRed] (-1.6,1.4) node[left, font=\sffamily\small, text=FpgaRed] {сброс} -- (s0);
\end{tikzpicture}
\caption{Граф переходов детектора \code{101} (автомат Мура). На стрелках ---
значение входа $x$}
\label{fig:fsm101}
\end{figure}

\begin{table}[H]
\centering
\caption{Таблица переходов детектора}
\begin{tabular}{l|cc|c}
\toprule
\textbf{Текущее состояние} & \multicolumn{2}{c|}{\textbf{Следующее состояние}} & \textbf{Выход} $y$ \\
 & $x=0$ & $x=1$ & \\
\midrule
S0 (ничего не найдено) & S0 & S1 & 0 \\
S1 (принято «1»)       & S2 & S1 & 0 \\
S2 (принято «10»)      & S0 & S3 & 0 \\
S3 (принято «101»)     & S2 & S1 & 1 \\
\bottomrule
\end{tabular}
\end{table}

\begin{figure}[H]
\centering
\begin{tikztimingtable}[timing/slope=0.05, xscale=1.3, yscale=1.4, timing/font=\sffamily\small, timing/rowdist=3.2ex]
  CLK     & 10{1L 1H} \\
  $x$     & 2H 2L 2H 2L 2H 2H 2L 2L 2H 2L \\
  state   & 1D{S0} 2D{S1} 2D{S2} 2D{S3} 2D{S2} 2D{S3} 2D{S1} 2D{S2} 2D{S0} 2D{S1} 1D{S2} \\
  $y$     & 5L 2H 2L 2H 9L \\
\end{tikztimingtable}
\caption{Работа детектора: вход \code{1010110010}. Выход равен 1 после
каждой найденной тройки \code{101}, в том числе перекрывающейся}
\end{figure}

Проверим по таблице: S0 --(1)$\to$ S1 --(0)$\to$ S2 --(1)$\to$ S3 ($y=1$)
--(0)$\to$ S2 --(1)$\to$ S3 ($y=1$) --(1)$\to$ S1 --(0)$\to$ S2 --(0)$\to$ S0
--(1)$\to$ S1.

\subsection{Автомат Мура на Verilog}
\label{sec:sync-fsm-v}

Структуру автомата (рис.~\ref{fig:fsm_struct}) принято переносить в код
\emph{тремя блоками}: регистр состояния, логика переходов, логика выходов.
Так по тексту сразу видно, где память, а где комбинационная логика.

\vfile{\codepath/b9_sync/fsm101.v}{code/b9\_sync/fsm101.v}

\textbf{Как читается.}
\begin{itemize}
  \item \code{localparam} даёт состояниям имена \code{S0}\ldots\code{S3} и
        двоичные коды; поменяв числа, можно перейти к другому кодированию
        (например, унитарному), не трогая остальной текст.
  \item Регистр \code{state} хранит текущее состояние, \code{next} --- то,
        которое будет после фронта. Первый блок (\code{posedge clk}) только
        копирует \code{next} в \code{state} (с синхронным сбросом в
        \code{S0}), и в нём стоит \code{<=}.
  \item Второй блок (\code{always @*}) --- таблица переходов, записанная
        оператором \code{case}: для каждого состояния и каждого значения
        \code{x}. Здесь стоит \code{=}, потому что это комбинационная логика.
        Ветка \code{default} страхует от неописанных кодов состояния. Здесь
        использованы все четыре кода, но при пяти состояниях и трёх битах без
        неё (и без значения \code{next} по умолчанию) вышла бы защёлка.
  \item Третий блок --- выход Мура: \code{y} равен 1 в состоянии \code{S3}
        и зависит только от состояния.
\end{itemize}

\textbf{Тот же автомат как автомат Мили.} Автомат Мили строят по тому же
шаблону, меняется только третий блок: выход зависит и от состояния, и от
входа. Состояний нужно три, четвёртое не требуется: вместо него «состояние
\code{S2} и на входе 1».

\vfile{\codepath/b9_sync/fsm101_mealy.v}{code/b9\_sync/fsm101\_mealy.v}

Тестбенч \code{tb\_fsm101.v} подаёт на оба автомата вход \code{1010110010} и
проверяет два факта: оба находят ровно две тройки \code{101} (с
перекрытием), а выход Мура запаздывает относительно выхода Мили ровно на
один такт.

\textbf{Как запустить в Logra.} Вкладка \term{Код}: вставьте \code{fsm101},
Ctrl+Enter. На вкладке \term{RTL} найдите регистр состояния (два триггера),
логику переходов и сравнение для выхода; на \term{Диаграммах} подайте сброс
на один такт, затем биты \code{1,0,1,0,1,1,0,0,1,0} по одному на такт и
сверьтесь с рисунком выше: \code{y} равен 1 после третьего и после пятого
бита. Для Мили повторите --- единица придёт на такт раньше.

\begin{diybox}[Автомат на Verilog: допиши модуль]
\textbf{Задание.} \stars{1} Измените \code{fsm101} так, чтобы он искал
\code{110}. \stars{2} Сделайте детектор \emph{без} перекрытия: после
найденной тройки автомат начинает искать с нуля. Чем отличаются таблицы
переходов? \stars{3} Перейдите к унитарному кодированию состояний:
четыре состояния в четырёх триггерах. Что изменилось на \term{RTL}?\\
\textbf{Ожидаемый результат:} для \code{110} единица появляется после
каждой тройки \code{1,1,0}; без перекрытия поток \code{10101} даёт одну
тройку, а с перекрытием две.
\end{diybox}

\textbf{Образец ответа} (детектор \code{110} с перекрытием, автомат Мура).
\begin{vcode}
module fsm110 (input clk, rst, x, output y);
    localparam S0 = 2'd0,    // ничего
               S1 = 2'd1,    // «1»
               S2 = 2'd2,    // «11»
               S3 = 2'd3;    // «110»
    reg [1:0] state, next;
    always @(posedge clk) begin
        if (rst) state <= S0;
        else     state <= next;
    end
    always @* begin
        case (state)
            S0:      next = x ? S1 : S0;
            S1:      next = x ? S2 : S0;
            S2:      next = x ? S2 : S3;   // «11» и снова 1: остаёмся в «11»
            S3:      next = x ? S1 : S0;   // после 110 новая тройка может начаться с 1
            default: next = S0;
        endcase
    end
    assign y = (state == S3);
endmodule
\end{vcode}

\subsection{Кодирование состояний}

Состояния хранятся в регистре в виде двоичных кодов. Четыре состояния можно
закодировать двумя битами (\code{00}, \code{01}, \code{10}, \code{11}) ---
это \term{двоичное кодирование}. Можно взять четыре триггера и держать единицу
только в одном из них (\code{0001}, \code{0010}, \ldots) --- это
\term{унитарное кодирование} (one-hot). Оно требует больше триггеров, зато
логика переходов проще и быстрее. Там, где триггеров много, обычно выбирают
именно его. \lookahead{как сделать выбор кодирования и посмотреть его цену в
числе элементов, мы увидим в главе~\ref{ch:fsmdesign} и в главе~\ref{ch:inside}.}

\begin{warnbox}[Типичные ошибки при работе с автоматами]
\begin{itemize}
  \item \textbf{Пустая ячейка таблицы переходов.} Для каждого состояния должна
        быть описана реакция на \emph{каждую} комбинацию входов, иначе
        автомат «зависнет» в неописанном случае.
  \item \textbf{Нет сброса.} Без начального состояния после включения
        регистр хранит случайный код; автомат может оказаться в состоянии, из
        которого нет выхода.
  \item \textbf{Путаница Мура и Мили.} Выход, который вы хотели получить в том
        же такте, у автомата Мура придёт на такт позже: он зависит только от
        состояния.
  \item \textbf{Одно \code{always} на всё.} Смешать состояние, переходы и
        выход в одном тактируемом блоке можно, но выход тогда вычисляется
        из \emph{следующего} состояния, и легко сдвинуть его на такт. Три
        блока (состояние, переходы, выход) этой ошибки избегают.
  \item \textbf{\code{=} вместо \code{<=} в регистре состояния} (и
        наоборот в блоке переходов): автомат работает в симуляторе иначе, чем
        после синтеза.
  \item \textbf{Нет \code{default} в \code{case} переходов.} Для неописанных
        кодов состояния \code{next} не присвоен, и синтезатор ставит защёлку.
  \item \textbf{Детектор без перекрытия.} В состоянии «найдено» нельзя
        возвращаться в начало: после \code{101} следующая единица уже может
        быть началом новой тройки.
\end{itemize}
\end{warnbox}

\begin{diybox}[Детектор и разновидности автоматов]
\textbf{Задание.} Собрать детектор последовательности, сравнить автоматы Мура
и Мили и проверить перекрытие.\\
\textbf{В Logra.} Вкладка \emph{Обучение}, книга «из блоков», глава
«Автоматы»; вкладки \emph{Конструктор} и \emph{Диаграммы}.\\
\stars{1} \emph{Повтори.} Урок «Детектор двух единиц подряд», затем урок
«Детектор 101».\\
\stars{2} \emph{Измени.} Уроки «Выход Мура» и «Выход Мили»: соберите оба для
одного условия и сравните диаграммы. Ожидаемо: выход Мили поднимается на такт
раньше.\\
\stars{3} \emph{Проект.} Проект \code{sequence-1011}: подайте поток
\code{1011011} после начальной синхронизации и убедитесь, что совпадение
находится дважды (перекрытие).\\
\textbf{Ожидаемый результат.} Автомат Мура выдаёт единицу в состоянии
«найдено», автомат Мили~--- в том же такте, где пришёл последний бит.\\
\textbf{Как проверить.} Вкладка \emph{Диаграммы}: сверьте последовательность
состояний с таблицей переходов главы; вкладка \emph{RTL}: найдите регистр
состояния и логику переходов.
\end{diybox}

\section{Углубление: память (ПЗУ и ОЗУ)}
\label{sec:memory}

\emph{Этот раздел можно прочитать позже: он не нужен для понимания автоматов,
но на нём строятся очередь данных следующей главы и процессор из главы~\ref{ch:cpu}.}

Регистр из $n$ триггеров хранит одно число. Чтобы хранить много чисел,
используют \term{память}: массив ячеек, каждая из которых имеет свой номер ---
\term{адрес}. Память характеризуется \term{глубиной} (числом ячеек) и
\term{шириной} (числом битов в ячейке), например $1024 \times 8$ ---
1024 ячейки по 8 бит.

\subsection{ПЗУ: память только для чтения}

\term{ПЗУ} (постоянное запоминающее устройство, ROM) хранит данные, записанные
заранее. По сути это комбинационная схема: адрес --- вход, данные --- выход,
а содержимое --- большая таблица истинности. ПЗУ можно построить из
дешифратора адреса и элементов ИЛИ: каждая ячейка --- это выход дешифратора,
подключённый к тем элементам ИЛИ, где в данных стоит 1.

\begin{figure}[H]
\centering
\begin{tikzpicture}
  \draw[thick, fill=FpgaLight] (0,-2.8) rectangle (1.4,0.6);
  \node[font=\sffamily\bfseries] at (0.7,-1.1) {DC};
  \draw[wire] (-0.7,-0.4) node[left] {$a_1$} -- (0,-0.4);
  \draw[wire] (-0.7,-1.8) node[left] {$a_0$} -- (0,-1.8);
  \foreach \i in {0,...,3} {
    \pgfmathsetmacro{\y}{0.2-\i*0.9}
    \draw[wire] (1.4,\y) -- (7.6,\y);
    \node[font=\sffamily\scriptsize, anchor=east] at (1.35,\y) {\i};
  }
  \foreach \j in {0,...,3} \draw[very thick, FpgaBlue] (2.4+\j*1.5,0.5) -- (2.4+\j*1.5,-3.2);
  % данные: ячейка0=1011, 1=0110, 2=1100, 3=0001 (d3..d0)
  \foreach \i/\row in {0/1011, 1/0110, 2/1100, 3/0001} {
    \pgfmathsetmacro{\y}{0.2-\i*0.9}
    \foreach \j in {0,...,3} {
      \StrChar{\row}{\numexpr\j+1\relax}[\bitv]
      \ifnum\bitv=1 \fill[FpgaRed] (2.4+\j*1.5,\y) circle (2.4pt);\fi
    }
    \node[font=\ttfamily\small, anchor=west] at (7.8,\y) {\row};
  }
  \foreach \j/\n in {0/3,1/2,2/1,3/0} {
    \node[gOR, rotate=-90] (o\j) at (2.4+\j*1.5,-3.8) {};
    \draw[wire] (o\j.east) -- ++(0,-0.4) node[below] {$d_\n$};
  }
  \node[lbl, anchor=west] at (7.8,0.75) {содержимое};
  \node[lbl, text=FpgaRed, align=left, anchor=west] at (7.8,-3.6) {точка --- соединение\\(в этом бите 1)};
\end{tikzpicture}
\caption{ПЗУ $4\times4$: дешифратор выбирает строку, элементы ИЛИ собирают
биты. \lookahead{тот же принцип --- «память как таблица истинности» --- лежит в основе таблиц LUT}}
\label{fig:rom}
\end{figure}

\subsubsection*{ПЗУ на Verilog}

На Verilog память --- это \emph{массив} слов: \code{reg [7:0] mem [0:15]}
означает 16 ячеек по 8 бит. Заполняют его в блоке \code{initial}, а читают
обычным индексом \code{mem[addr]}.

\vfile{\codepath/b9_sync/rom16x8.v}{code/b9\_sync/rom16x8.v}

\textbf{Как читается.} В блоке \code{initial} каждой ячейке присвоено слово:
это коды сегментов семисегментного индикатора для цифр \code{0}\ldots\code{F}
(бит 0 --- сегмент $a$, бит 6 --- сегмент $g$). Строка
\code{assign data = mem[addr]} --- комбинационное чтение: адрес выбирает
слово, как в схеме на рис.~\ref{fig:rom}. Обновления по такту здесь нет:
значение на выходе меняется сразу после смены адреса. Синтезатор превращает
такой массив в таблицу истинности (или в блок памяти ПЛИС).

\subsection{ОЗУ: оперативная память}

\term{ОЗУ} (оперативное запоминающее устройство, RAM) позволяет и читать, и
записывать. Концептуально это набор регистров, дешифратор адреса, который при
записи выбирает, в какой регистр записать данные, и мультиплексор, который
при чтении выбирает, какой регистр подать на выход.

\begin{figure}[H]
\centering
\begin{tikzpicture}
  \draw[thick, fill=FpgaLight] (0,-2.9) rectangle (1.3,0.7);
  \node[font=\sffamily\bfseries] at (0.65,-1.1) {DC};
  \foreach \i in {0,...,3} {
    \pgfmathsetmacro{\y}{0.25-\i*0.95}
    \node[draw=FpgaGreen, fill=green!8, thick, minimum width=22mm, minimum height=7mm, font=\sffamily\small] (r\i) at (3.6,\y) {ячейка \i};
    \draw[arr] (1.3,\y) -- node[above, lbl] {CE} (r\i.west);
    \draw[bus, -{Stealth}] (r\i.east) -- (5.9,\y);
  }
  \draw[thick, fill=orange!15] (5.9,0.7) -- (7.0,0.2) -- (7.0,-2.4) -- (5.9,-2.9) -- cycle;
  \node[font=\sffamily\small] at (6.45,-1.1) {MUX};
  \draw[bus, -{Stealth}] (7.0,-1.1) -- (8.2,-1.1) node[right, text=FpgaDark] {данные чтения};
  \draw[bus] (3.6,1.6) node[above, text=FpgaDark] {данные записи} -- (3.6,0.6);
  \draw[arr] (-0.8,0.3) node[left, font=\sffamily\small] {WE (запись)} -- (0,0.3);
  \draw[bus, -{Stealth}] (-0.8,-1.1) node[left, text=FpgaDark] {адрес} -- (0,-1.1);
  \draw[bus] (-0.4,-1.1) |- (6.45,-3.6) -- (6.45,-2.65); \fill[FpgaBlue] (-0.4,-1.1) circle (1.8pt);
\end{tikzpicture}
\caption{Устройство ОЗУ «на пальцах»: дешифратор адреса при записи,
мультиплексор при чтении}
\label{fig:ram}
\end{figure}

На практике ячейки памяти делают не из D-триггеров, а из более компактных
ячеек:
\begin{itemize}
  \item \textbf{статическая память SRAM} --- 6 транзисторов на бит, быстрая,
        не требует обслуживания; из неё сделаны кэши процессоров и блоки
        памяти внутри микросхем;
  \item \textbf{динамическая память DRAM} --- 1 транзистор и 1 конденсатор на
        бит, очень плотная, но требует периодической регенерации; именно
        такая память стоит в компьютерах как основная.
\end{itemize}

\subsubsection*{ОЗУ на Verilog}

Тот же массив, но с записью по такту:

\vfile{\codepath/b9_sync/ram16x8.v}{code/b9\_sync/ram16x8.v}

\textbf{Как читается.} По фронту, если \code{we}, слово \code{din}
записывается в ячейку с номером \code{addr}. Следующая строка читает ячейку в
регистр \code{dout}. Оба присваивания неблокирующие, поэтому чтение видит
содержимое \emph{до} фронта: если в этом же такте шла запись, на выходе
окажется старое слово, а новое --- на следующем такте. Так устроены блоки
памяти ПЛИС, и поэтому «чтение занимает такт».

Тестбенч \code{tb\_mem.v} проверяет ПЗУ (слова 0, 9 и 15), записывает четыре
слова в ОЗУ и читает их по одному, а затем проверяет, что чтение во время
записи отдаёт старое слово.

\textbf{Как запустить в Logra.} Вкладка \term{Код}: вставьте \code{rom16x8},
Ctrl+Enter, \term{Диаграммы}: подайте адреса от 0 до 15 и сверьте слова
с таблицей кодов в тексте модуля. Для \code{ram16x8}: подайте
\code{we=1}, адрес 3, \code{din=0x5A} на один такт, затем \code{we=0},
адрес 3 --- слово \code{5A} появится на \code{dout} через такт.

\begin{diybox}[ПЗУ и ОЗУ своими руками]
\textbf{Задание.} Прочитать заранее записанное слово по адресу, затем
записать слово и прочитать его обратно.\\
\textbf{В Logra.} Вкладка \emph{Примеры}: \code{small-rom},
\code{rom-melody}, \code{ram-write-read}; вкладка \emph{Обучение}, книга по
Verilog, главы про память: уроки «ПЗУ из массива» и «ОЗУ 8×4: запись по
такту».\\
\stars{1} \emph{Повтори.} Проект \code{small-rom}: выберите по очереди все
восемь адресов и выпишите слова; урок «ПЗУ из массива».\\
\stars{2} \emph{Измени.} В \code{small-rom} замените слово одного адреса и
проверьте, что соседние адреса не изменились; в \code{rom-melody} переставьте
два слова таблицы и сократите \code{STEP\_TICKS}.\\
\stars{3} \emph{Проект.} Проект \code{ram-write-read}: запишите четыре
слова и считайте их; сопоставьте каждое прочитанное слово с предыдущим
адресом (чтение занимает такт).\\
\textbf{Ожидаемый результат.} ПЗУ при каждом чтении выдаёт одно и то же,
ОЗУ запоминает записанное, пока не будет записано новое.\\
\textbf{Как проверить.} Вкладка \emph{Диаграммы}: адрес, слово записи,
сигнал разрешения записи и слово чтения; вкладка \emph{RTL}: сравните,
как выглядят ПЗУ (выбор готового слова по адресу) и ОЗУ (ячейки, которые
запоминают).
\end{diybox}

\textbf{Допиши модуль.} \stars{1} Измените слово по адресу 3 в
\code{rom16x8} (например, на \code{8'hFF}): остальные адреса должны остаться
прежними. \stars{2} Сделайте ПЗУ $8\times7$ с цифрами 0--7 и счётчиком
адреса, который перебирает слова по такту (мини-проект из задач к главе).
\stars{3} Добавьте к ОЗУ сброс содержимого в 0 по входу \code{clr}: сколько
тактов потребуется? \emph{Образец к заданию 2:}

\vfile{\codepath/b9_sync/seg_scan.v}{code/b9\_sync/seg\_scan.v}

Счётчик \code{addr} на три бита сам возвращается из 7 в 0, ПЗУ выдаёт слово
по текущему адресу; тестбенч \code{tb\_seg\_scan.v} проверяет 16 тактов подряд:
выход проходит все 8 слов в порядке адресов дважды. В задании~3 за один
такт можно очистить только одну ячейку (у блока памяти один порт записи),
поэтому при \code{clr=1} счётчик адреса перебирает ячейки и очистка занимает
16 тактов.

\section{Вопросы и задачи}
\label{sec:tasks-sync}

Ответы собраны в приложении~\ref{app:answers}.

\begin{taskbox}[Задачи к главе]
\begin{enumerate}
  \item \stars{1}~Максимальная задержка комбинационной логики между регистрами 8~нс, время установки 0,5~нс, задержка выхода триггера 0,5~нс. Чему равна максимальная тактовая частота?
  \item \stars{2}~Как изменится период такта, если разбить ту же логику на два одинаковых этапа конвейера? Что при этом происходит с задержкой результата?
  \item \stars{1}~Чем автомат Мура отличается от автомата Мили? К какому типу относится детектор \code{101} из этой главы (рис.~\ref{fig:fsm101}) и как изменился бы его граф, если бы выход $y$ нужно было получить в том же такте, когда пришёл последний бит тройки?
  \item \stars{1}~Сколько состояний нужно детектору последовательности 101 (с перекрытием)? Сколько триггеров потребуется при двоичном кодировании?
  \item \stars{1}~ПЗУ организовано как $256\times 8$. Сколько у него адресных линий и чему равна ёмкость в битах?
  \item \stars{1}~ОЗУ ёмкостью $1\text{К}\times 16$ бит: сколько это килобайт? Сколько адресных и информационных линий?
  \item \stars{2}~\emph{Найдите ошибку.} Студент рисует детектор \code{101} с перекрытием и в состоянии S3 при входе 1 возвращает автомат в S0. На входе \code{101101} такой автомат выдаст не то, что нужно. Сколько единиц выхода $y$ получится вместо правильных двух и в каком такте проявится ошибка?
  \item \stars{3}~\emph{Мини-проект.} Постройте в Конструкторе (или в Коде) детектор последовательности \code{1101} с перекрытием. Критерий приёмки: на диаграмме для входа \code{11011011101} выход равен 1 ровно столько раз, сколько раз в потоке встречается \code{1101} (подсчитайте заранее); таблица переходов без пустых ячеек.
  \item \stars{3}~\emph{Мини-проект.} Опишите в Коде ПЗУ $8\times 7$, которое хранит коды сегментов для цифр 0--7 семисегментного индикатора, и счётчик адреса, который перебирает её по такту. Критерий приёмки: на диаграмме за 8 тактов выход проходит все 8 слов в порядке адресов, а при замене одного слова меняется ровно один такт.
\end{enumerate}
\end{taskbox}

\section{Что дальше}

Автомат из этой главы мы разобрали на одном примере. Но как спроектировать
автомат для задачи, которой нет в учебнике: как не пропустить случай, как
выбрать состояния, как проверить решение? В главе~\ref{ch:fsmdesign} мы
превратим пример в рецепт из шести шагов, разберём автомат по продаже и
построим очередь данных~--- на той же памяти, что описана в разделе
«Углубление» этой главы. Сквозной проект «Охранная панель» там получит свой
главный этап: автомат режимов.
